\scanpage{72}
\mspage{16}

I) \uline{Définition de} $\underline{\Phi},\Psi$.

On va calquer la construction \inserted{canonique} de $\underline{\mathrm{Ch}}(X)\to X\times X$, \inserted{(ou $\underline{\mathrm{Ch}}_{\unsure{\infty|w}}(X)\unsure{\text{---}|\to}X\times X$\unsure{,})} factorisant $X\xrightarrow{\mathrm{diag}}X\times X$, et la \deleted{\unsure{constr.}} factorisation canonique correspondante d'un $f:X'\unsure{\text{---}|\to}X$ en $X'\xrightarrow{\deleted{\unsure{\sim}}\,i}\overline{X}'\xrightarrow{\bar f}Y$. On remplacera simplement $\underline{\mathrm{Ch}}(X)$ par la sous-catégorie $\underline{\mathrm{Ch}}_{\deleted{\unsure{\Delta}}+}(X)$, laquelle n'est autre que la catégorie des flèches de $X$
\[ \underline{\mathrm{Fl}}(X)=\deleted{\unsure{\mathrm{Hom}}}\,\underline{\mathrm{Hom}}(\Delta_1,X), \]
muni des foncteurs
\[ X\xrightarrow{i_X}\underline{\mathrm{Fl}}(X)\xrightarrow{\inserted{p_X=(\sigma_X,\beta_X)}} X\times X, \]
où $i_X$ est la flèche $x\longmapsto\mathrm{id}_x$, et \deleted{$\unsure{i_X(x)=}$} $\sigma_X,\beta_X$ sont les foncteurs source et but. La fibre de $p_X$ \inserted{en $(a,b)\in\operatorname{Ob}(X\times X)$} est la catégorie discrète définie par l'ens $\mathrm{Hom}(a,b)$, \unsure{ça} dépend \unsure{grandement} de $(a,b)$ et pas seulement de sa composante connexe dans $X\times X$ --- donc $p_X$ n'a rien d'une Hot-fibration. Il n'est en général ni propre ni lisse\unsure{,} ni rien, par contre les deux flèches composées avec $\mathrm{pr}_1$ et $\mathrm{pr}_2$
\diagram{\begin{tikzcd}[ampersand replacement=\&,column sep=small]
\&\underline{\mathrm{Fl}}(X)\arrow[dl,"\sigma_X"']\arrow[dr,"\beta_X"]\&\\
X\&\& X
\end{tikzcd}}
sont sympathiques. on a en effet

% archive sheet number: 71
